\chapter{Практика}


\section{Ход работы}

\section{Работа с терминалом Ubuntu}

\subsection{sudo}
Sudo (Substitute User and Do) переводится как «заменить пользователя и выполнить».
Данная команда позволяет нам выполнять команды с
привилегиями другого пользователя (суперпользователя, как правило).

\begin{figure}[H]
    \centering
    \includegraphics{ub1.png}
    \caption{Пример работы sudo}
\end{figure}

\subsection{man}
Man (от английского «manual») в Linux и других подобных ОС используется для доступа к
справочной информации по какой-либо команде, утилите или функции системы. 
У каждой команды есть свой синтаксис, в котором нужно разбираться. Тут
нам и может помочь команда man. 

\begin{figure}[H]
    \centering
    \includegraphics{ub2.png}
    \caption{Пример работы man}
\end{figure}

\subsection{Полезные сочетания клавиш}
\begin{enumerate}
    \item Ctrl+c отправляет команду текущему процессу на его немедленное завершение.
    \item Ctrl+z отправляет команду текущему процессу на его приостановку
\end{enumerate}


\section{Работа с командной строкой Windows}

\subsection{RunAs}
Как таковой команды-аналога sudo в системах Windows нет (её вводят только в новых
версиях Windows 11). Для получения привилегированного доступа можно воспользоваться
командой runas. Смысл этой команды такой же, но она потребует от вас ввода полных данных
пользователя с правами администратора

\begin{figure}[H]
    \centering
    \includegraphics{win1.png}
    \caption{Пример работы runas}
\end{figure}

\subsection{/?}
Также, как и с командой man, справочную информацию в системах Windows можно
получить через команду /?, используя данный синтаксис после любой команды. Например,
ping /?. И мы получаем вывод справочной информации по данной команде.

\begin{figure}[H]
    \centering
    \includegraphics{win2.png}
    \caption{Пример работы /?}
\end{figure}

\subsection{Полезные сочетания клавиш}
\begin{enumerate}
    \item Ctrl+c отправляет команду текущему процессу на его немедленное завершение.
    \item Ctrl+s отправляет команду текущему процессу на его приостановку
\end{enumerate}

\endinput